\hypertarget{testing-and-verification}{%
\chapter{Testing and Verification}\label{testing-and-verification}}

\begin{draftnote}

\textbf{Draft status, 04 November 2026.} Figures read from
\texttt{docs/TESTING.md}, \texttt{apps/api/package.json},
\texttt{apps/api/tests/helpers.ts}, \texttt{.github/workflows/ci.yml}
and \texttt{docs/PROJECT\_NOTES.md}. Re-run the suite and confirm the
counts before submission; a test count quoted wrongly is the easiest
thing for an examiner to check.

\end{draftnote}

\begin{center}\rule{0.5\linewidth}{0.5pt}\end{center}

\hypertarget{what-is-tested-and-what-that-number-means}{%
\section{What is tested, and what that number
means}\label{what-is-tested-and-what-that-number-means}}

\textbf{278 automated tests: 155 in the API under vitest, 123 in the
detector under pytest.}

The split is not arbitrary. The detector suite covers the whole
structural tier --- parsing, normalisation, tree edit distance, the
prefilter, unit matching, the style-feature extraction and both HTTP
endpoints --- and needs no database, no object store and no Docker. It
runs in under three seconds and is the fastest useful signal in the
project. The API suite covers authentication, courses, submissions,
results, the configuration resolver, Layer 2 and the re-run path, and
needs MongoDB, Redis and the object store running.

A count on its own is close to meaningless, so the rest of this chapter
is about what the suite is constructed to prevent rather than how much
of it there is.

\hypertarget{three-design-decisions-in-the-test-harness}{%
\section{Three design decisions in the test
harness}\label{three-design-decisions-in-the-test-harness}}

These are worth a chapter section because each one exists in response to
a specific way the suite could have produced false results.

\hypertarget{a-guard-that-refuses-to-run-against-the-development-database}{%
\subsection{A guard that refuses to run against the development
database}\label{a-guard-that-refuses-to-run-against-the-development-database}}

\texttt{tests/helpers.ts} rewrites the Mongo URI to point at
\texttt{codeguard\_test}, and then refuses to proceed if the rewrite did
not take effect:

\begin{verbatim}
process.env.MONGO_URI = (process.env.MONGO_URI ?? "").replace("/codeguard?", `/${TEST_DB}?`);
if (!process.env.MONGO_URI.includes(`/${TEST_DB}?`)) {
  throw new Error(`Refusing to run: tests must use the ${TEST_DB} database`);
}
\end{verbatim}

The suite clears collections between tests. Without this guard, a
mistyped or unset URI would point that clearing at the development
database. The guard is the thing standing between a typo and the loss of
every fixture, course and submission in local development.

It has a consequence worth knowing: the \texttt{.env} path is
\textbf{relative to the working directory}, so running the suite from
the repository root loads nothing, leaves \texttt{MONGO\_URI} empty, and
produces 155 identical failures on that error message. \textbf{That is
the guard working, not 155 real failures} --- a distinction that costs a
quarter of an hour the first time it happens.

\hypertarget{serial-execution-because-the-files-share-one-database}{%
\subsection{Serial execution, because the files share one
database}\label{serial-execution-because-the-files-share-one-database}}

\begin{verbatim}
"test": "vitest run --no-file-parallelism"
\end{verbatim}

Every test file connects to the same database and clears collections
between tests. Run the files in parallel and they delete each other's
fixtures mid-run. The failures that produces are real-looking, entirely
fictional, and move around between runs --- the worst category of red to
debug, because it punishes the correct instinct of re-running to
confirm.

\texttt{npm\ test} also runs \texttt{pretest}, which typechecks first,
so a drifted contract is caught in seconds before any container is
touched. Calling \texttt{vitest} directly loses both properties, which
is why the project's documentation specifies \texttt{npm\ test} rather
than the tool underneath it.

\hypertarget{the-detector-is-stubbed-deliberately}{%
\subsection{The detector is stubbed,
deliberately}\label{the-detector-is-stubbed-deliberately}}

\texttt{helpers.ts} overrides \texttt{DETECTOR\_URL} to a stub port
\emph{after} loading the environment, with the comment ``so the real
detector can never be reached''. A detector that is down, mid-rebuild or
on a different version therefore cannot make the API suite red.

This is a deliberate trade. The API suite tests the API's own behaviour,
including how it handles what the detector returns, and isolating it
means a red API suite always indicates an API defect. The cost is that
no automated test exercises the two services together; that gap is
stated in \S{}6.6.

\hypertarget{mutation-testing}{%
\section{Mutation testing}\label{mutation-testing}}

This is the strongest verification evidence in the project and the part
most likely to be asked about.

A passing suite demonstrates that the tests agree with the code. It does
not demonstrate that the tests would notice if the code were wrong.
Mutation testing asks the second question directly: deliberately break
the implementation in a small, plausible way, and check that at least
one test fails.

\textbf{Nineteen deliberate mutations were introduced across the
detector suite. Seventeen were caught immediately. Two were not}, and
those two are the valuable ones:

\begin{itemize}
\tightlist
\item
  One survived because it mutated a code path \textbf{no test reached}
  --- a dead branch the suite had never exercised and the coverage
  figure had not revealed as untested in any way that prompted action.
\item
  One survived because the suite asserted the \emph{contents} of a
  result but never its \textbf{ordering}, so reversing a sort changed
  nothing any assertion looked at.
\end{itemize}

Both gaps were then closed. The finding to report is not ``19 of 19
eventually caught'' but that \textbf{a suite can be green, substantial
and still blind in a specific direction}, and that mutation testing is
what makes that blindness visible. Very few undergraduate projects
contain this, and it directly answers the question \emph{how do you know
your tests are any good?}

\hypertarget{continuous-integration-and-a-gate-that-was-not-shut}{%
\section{Continuous integration, and a gate that was not
shut}\label{continuous-integration-and-a-gate-that-was-not-shut}}

The CI workflow runs on every push and pull request against
\texttt{main}, with four jobs: API typecheck and production-dependency
audit, \textbf{API tests}, detector lint and tests, and web typecheck.

\textbf{The API tests job did not exist until 01 November 2026, and the
honest account of that belongs in the report.} For most of the project
the workflow had three jobs --- API typecheck, detector lint and tests,
web typecheck --- and none of them ran vitest. The gate therefore
covered 123 of 278 tests, and every green badge meant \emph{the types
check and the detector passes} while being read as \emph{the tests
pass}.

A badge that means less than it appears to is worse than no badge,
because it is trusted. The gap was found while investigating an
unrelated red badge, which is its own small lesson: the failure was
interesting, and what it revealed was more interesting.

Closing it required MongoDB and Redis as service containers with health
checks, and the object store started as a \texttt{docker\ run} step
rather than a service --- because an Actions service container cannot be
given a command, and the object store needs one. The stated reason for
not doing it earlier had been that it would require putting credentials
into GitHub secrets; two thirds of that reasoning was wrong.
\texttt{helpers.ts} loads the environment file only
\texttt{if\ (existsSync(...))}, so a missing file is not an error and
the workflow can supply the environment itself, and because every
container is destroyed with the runner, nothing it holds is a credential
worth protecting.

\hypertarget{the-dependency-audit-as-a-gate}{%
\subsection{The dependency audit as a
gate}\label{the-dependency-audit-as-a-gate}}

\texttt{npm\ run\ audit:prod} fails the build on a high or critical
advisory against a production dependency. It is advisory-aware: a small
allow-list records specific advisory identifiers with a written reason
and a date, so a known and assessed issue does not block every
subsequent push while an unassessed one still does.

One episode is worth recording because it is a general lesson about
fixing advisories. A transitive dependency had a published advisory
whose patched versions were all ESM-only, while the consuming package is
CommonJS and calls the export directly. Applying the ``fix'' would have
replaced a working function export with a module object and broken the
object-store client at runtime --- a security patch that would have
taken down file uploads. \textbf{The advisory was real, the patch was
real, and applying it would have been wrong.} It was caught by comparing
the exported shapes rather than by trusting the version number.

\hypertarget{a-verification-method-the-suite-does-not-cover}{%
\section{A verification method the suite does not
cover}\label{a-verification-method-the-suite-does-not-cover}}

Three classes of defect in this project were found by methods other than
running tests, and the counts are worth stating because they argue for
keeping those methods rather than replacing them with more tests.

\textbf{Rendering real data through a screen and reading it.} Nine
defects were found this way across two days of frontend work, and more
later. They included a CSS custom property that did not exist --- a bad
\texttt{var()} fails silently, so the colour simply came out wrong --- a
label describing a cohort statistic as belonging to the student, and two
distinct meanings sharing one visual language.

\textbf{Printing the result to paper.} Two further defects, including a
background wash that appeared on every printed page and came from a
container whose padding had been reset for print but whose background
had not.

\textbf{Reading the code rather than executing it.} Four defects in the
ablation study, described in Chapter 7 \S{}7.2. Each had produced output
that looked like a result: a positive control that compared the wrong
file, a chance baseline that inflated every figure, a measurement of a
comparison strategy the system does not deploy, and a significance
verdict printed without its direction.

What these have in common is that \textbf{a test asserts a value, and
every one of these was a thing nobody had thought to assert.} They are
not an argument against automated testing; they are an argument that a
suite's coverage is bounded by its author's imagination, and that
rendering, printing and reading are the cheap ways to exceed it.

\hypertarget{limitations-of-the-verification}{%
\section{Limitations of the
verification}\label{limitations-of-the-verification}}

Stated here rather than left for an examiner to find.

\begin{enumerate}
\def\labelenumi{\arabic{enumi}.}
\tightlist
\item
  \textbf{No end-to-end test exercises the API and detector together.}
  The detector is stubbed in the API suite by design (\S{}6.2.3), so the
  integration between the two services is verified manually and by the
  seeded demonstration script.
\item
  \textbf{No automated frontend tests.} The web interface is verified by
  rendering and inspection, which \S{}6.5 shows is effective at finding a
  particular class of defect and is not a substitute for regression
  tests.
\item
  \textbf{Mutation testing covers the detector only.} The API suite has
  not been subjected to it, so the equivalent blind spots there are
  unknown rather than absent.
\item
  \textbf{The API suite takes several minutes}, dominated by password
  hashing. One test signs in eight times at roughly half a second of
  bcrypt each, which left it one slow machine away from a spurious
  timeout until it was given its own budget. Slow suites get run less
  often, and that is a verification risk rather than an inconvenience.
\item
  \textbf{Coverage percentage is not reported.} It was not treated as a
  goal, on the grounds that mutation testing answers the question
  coverage is usually used as a proxy for. An examiner may still ask for
  the number.
\end{enumerate}
